\section{Internal moments and scales of compactness}
\label{tree:sec:moments}

We seek deterministic scales at which internal graph distances are tight
and their normalizing moment survives in the limit. The ingredients are
the unconditioned bilateral walk of
Theorem~\ref{tree:thm:background_mullin}, independence of its increments,
and the volume bound \eqref{tree:eq:background_volume}.

For an integer interval $I=[a,b]$ of length $|I|=b-a$, let $G_I$,
$d_I$, and $\pi_I$ be its internal graph, all-edge distance, and contour
projection. Define the endpoint costs, diameters, and dyadic endpoint norms by
\begin{equation}\label{tree:eq:internal_norms}
 X_I=d_I(\pi_I(a),\pi_I(b)),\qquad
 \Delta_I=\diam(G_I,d_I),\qquad
 A_j=\bigl(\E X_{[0,2^j]}^k\bigr)^{1/k}\quad(j\geq0).
\end{equation}
Graph diameters refer to vertex sets. Recall that $p=1/d_\gamma$,
$s=p/2$, and the integer $k$ is fixed with $ks>8$; in particular,
$s>1/k$.

The contour path gives $X_I\leq |I|$. Equality holds when all increments
are $(1,0)$, an event of positive probability, so every $A_j$ is finite
and positive; in fact $A_0=2^{-1/k}$. Internal paths remain available
when their interval is enlarged. Concatenation and Minkowski's inequality
therefore give
\begin{equation}\label{tree:eq:internal_subadditivity}
 X_{[a,c]}\leq X_{[a,b]}+X_{[b,c]}
 \quad(a\leq b\leq c),\qquad A_{j+1}\leq2A_j.
\end{equation}
Costs on intervals with disjoint interiors are independent, since each
uses only its own increments. Their laws depend only on the interval
lengths.

\Needspace{10\baselineskip}
\subsection{Norm growth from volume}

The volume-counting idea is closely related to
\cite[Lemma~3.2]{GwynnePfefferDLA}: many vertices visited during an
encoded interval must fit in one ambient metric ball. We apply it below
to endpoint moments and prove the needed estimate directly.

\begin{lemma}[Norm growth]\label{tree:lem:moment_growth}
The deterministic sequence in \eqref{tree:eq:internal_norms} satisfies
\begin{equation}\label{tree:eq:moment_limsup}
 \limsup_{j\to\infty}\frac{\log_2 A_j}{j}\geq p.
\end{equation}
\end{lemma}

\begin{proof}
For $0\leq r\leq j$, let $M_{j,r}$ be the maximum of $X_J$ over the
$2^r$ consecutive dyadic intervals $J$ of length $2^{j-r}$ in
$[0,2^j]$. For nonnegative random variables, the $k$th power of their
maximum is at most the sum of their $k$th powers. Thus
\begin{equation}\label{tree:eq:dyadic_max_basic}
 \|M_{j,r}\|_{L^k}\leq2^{r/k}A_{j-r}.
\end{equation}
Decompose $[0,u]$, for any integer $u\in[0,2^j]$, into dyadic intervals
with at most one of each length. Their internal endpoint paths concatenate
to an all-edge path in the full map. Thus, with $\pi$ denoting the
full-map projection,
\[
 \max_{0\leq u\leq2^j}d_{\mathcal M}(\pi(0),\pi(u))
 \leq\sum_{r=0}^jM_{j,r}.
\]
Suppose \eqref{tree:eq:moment_limsup} fails. Since $1/k<p$, we may choose
$1/k<\alpha<\beta<p$ and a finite $C$ such that
$A_j\leq C2^{\alpha j}$ for every $j$. By Minkowski's inequality and
\eqref{tree:eq:dyadic_max_basic}, the $L^k$ norm of the last display's
right-hand side is at most
\[
 C2^{\alpha j}\sum_{r=0}^j2^{-r(\alpha-1/k)}
 \leq C'2^{\alpha j}.
\]
Since $\alpha<\beta$, Markov's inequality places every vertex visited
during $[0,2^j]$ in the ambient ball of radius $\lceil2^{\beta j}\rceil$
about $\pi(0)$, with probability tending to one.

The interval has $2^j/4+o_{\Prob}(2^j)$ first-coordinate up-steps, whose
terminal vertices are distinct. Different terminal heights cannot be
equivalent. If two terminal heights both equal $h$, the base of the later
step has height $h-1$ between the two terminal times, violating the minimum
criterion for equivalence. This also holds in the bilateral map. The ball
therefore contains at least $2^j/8$ vertices with probability tending to
one. Choose $\varepsilon>0$ with $\beta(d_\gamma+\varepsilon)<1$.
By \eqref{tree:eq:background_volume}, the same ball has at most
$\lceil2^{\beta j}\rceil^{d_\gamma+\varepsilon}=o(2^j)$ vertices with
probability tending to one, a contradiction.
\end{proof}

\subsection{A quantitative dyadic modulus}

A bound on earlier endpoint norms controls all internal distances and
their time modulus. We keep the denominator arbitrary: the estimate
will be needed before the limiting norm at a larger scale is known to
be positive.

\begin{lemma}[Moment compactness]\label{tree:lem:moment_compactness}
Fix $K<\infty$. Suppose $j\geq0$ and a deterministic $B>0$ satisfy
\begin{equation}\label{tree:eq:past_ratio_bound}
 \frac{A_{j-r}}{B}\leq K2^{-sr},\qquad 0\leq r\leq j.
\end{equation}
For each $T\geq1$ there is a finite $C_T$, depending only on
$T,k,s$, such that every deterministic integer interval $I$ of length
$1\leq m\leq T2^j$ satisfies
\begin{equation}\label{tree:eq:internal_diameter_bound}
 \left\|\frac{\Delta_I}{B}\right\|_{L^k}
 \leq C_TK\left(\frac{m}{2^j}\right)^s.
\end{equation}
Put $q=s-1/k>0$. For $2^{-j}\leq\delta\leq1$, set
\[
 \Omega_{j,T}(\delta)=\frac1B
 \max\left\{X_{[a,b]}:
  a,b\in\Z\cap[-T2^j,T2^j],\ 0\leq b-a\leq\delta2^j\right\}.
\]
Then
\begin{equation}\label{tree:eq:uniform_time_modulus}
 \|\Omega_{j,T}(\delta)\|_{L^k}\leq C_TK\delta^q.
\end{equation}
In particular, along any sequence $j\to\infty$ satisfying
\eqref{tree:eq:past_ratio_bound} with a common $K$, $B\to\infty$;
the rescaled pullbacks of $d_J$, for a fixed bounded macroscopic
confinement window $J$, are tight after interpolation and every
limit is a continuous pseudometric. The same estimates apply on every fixed finite
enlargement of the base window.
\end{lemma}

\begin{proof}
Every integer interval $[a,b]$ is the union of its maximal contained
dyadic intervals, with at most two of each length. Their endpoint paths
concatenate inside $[a,b]$ and hence are available in every $G_I$ with
$[a,b]\subset I$. We use the dyadic grid on the full integer line.

\emph{Diameter bound.} First assume $m\leq2^j$ and choose
$r_0\in\{0,\ldots,j\}$ so that
$2^{j-r_0}\leq m<2^{j-r_0+1}$. At depth $r\geq r_0$, let $M_r(I)$
be the maximum of $X_J/B$ over dyadic intervals $J\subset I$ of
length $2^{j-r}$, with the maximum of the empty collection set to
zero. There are at most $2^{r-r_0+1}$ such intervals. Consequently
\[
 \|M_r(I)\|_{L^k}
 \leq 2^{(r-r_0+1)/k}K2^{-sr}.
\]
For any two vertices of $G_I$, choose contour representatives
$a\leq b$ and use the preceding dyadic decomposition. It follows
simultaneously for all pairs that
\[
 \frac{\Delta_I}{B}\leq2\sum_{r=r_0}^j M_r(I).
\]
The geometric series gives an upper bound
$2^{1+1/k}K2^{-sr_0}/(1-2^{-q})$ for its $L^k$ norm. Since
$2^{-r_0}\leq m/2^j$, this proves
\eqref{tree:eq:internal_diameter_bound} for $m\leq2^j$. When
$2^j<m\leq T2^j$, split $I$ into at most $\lceil T\rceil$ pieces
of length at most $2^j$. Concatenation and the first bound prove
the assertion after increasing $C_T$.

\emph{Time modulus.} Choose $r_0$ so that
$2^{-r_0}\leq\delta<2^{-r_0+1}$. At depth $r\geq r_0$, there are
at most $(2T+2)2^r$ dyadic intervals of length $2^{j-r}$ in the
window. The $L^k$ norm of their maximum normalized endpoint cost is
at most $(2T+2)^{1/k}K2^{-qr}$. Every interval of length at most
$\delta2^j$ has at most two dyadic pieces at each such depth, so
summing this bound proves \eqref{tree:eq:uniform_time_modulus}.

\emph{Compactness.} Taking $r=j$ in \eqref{tree:eq:past_ratio_bound} gives
$B\geq A_0K^{-1}2^{sj}$, hence $B\to\infty$. On a fixed window
$J$, the triangle inequality bounds the change in $d_J(a,b)/B$
when its arguments change by the corresponding two short-interval
costs. Equation~\eqref{tree:eq:uniform_time_modulus}, with the harmless
replacement of $\delta$ by $\delta+2^{-j}$ for rounding, controls
this change. The diameter bound gives tight values. The
Arzel\`a--Ascoli criterion in probability therefore gives tightness
of continuous interpolations, whose interpolation errors vanish.
Symmetry, nonnegativity, zero diagonal, and the triangle inequality
pass to every limit. The estimates used the full translated grid,
so they hold on bounded enlargements as well.
\end{proof}

\subsection{Preserving the normalizing moment}

Moment compactness alone does not ensure that the normalizing moment
passes to the limit. Independence of disjoint walk intervals supplies
the additional uniform integrability.

\begin{lemma}[Uniform integrability without a higher moment]
\label{tree:lem:moment_ui}
Let $j_n\to\infty$ and deterministic $B_n>0$ satisfy
\eqref{tree:eq:past_ratio_bound} with a common $K$. Fix $T<\infty$.
For every sequence of deterministic integer intervals $I_n$ of length at most
$T2^{j_n}$, the variables $(X_{I_n}/B_n)^k$ are uniformly
integrable. Consequently, if
$X_{I_n}/B_n\Rightarrow X$, then
\begin{equation}\label{tree:eq:moment_convergence}
 \E X^k=\lim_{n\to\infty}\E(X_{I_n}/B_n)^k.
\end{equation}
This conclusion also holds for any nonnegative variables bounded
by the corresponding $X_{I_n}/B_n$, including distances between
the same endpoints computed in a larger internal graph. In
particular, when $B_n=A_{j_n}$ and $|I_n|=2^{j_n}$, any weak
limit of $X_{I_n}/A_{j_n}$ has $k$th moment one.
\end{lemma}

\begin{proof}
Set $X_n=X_{I_n}/B_n$. Fix an integer $m\geq2$ and split $I_n$
into $m$ consecutive integer intervals with lengths differing by
at most one; zero-length pieces are allowed. Let $Y_{n,1},\ldots,
Y_{n,m}$ be their normalized internal endpoint costs. They are
independent for fixed $n$, and $X_n\leq S_n:=\sum_iY_{n,i}$.
By Lemma~\ref{tree:lem:moment_compactness}, for fixed $m$ and all
sufficiently large $n$,
\begin{equation}\label{tree:eq:ui_piece_bound}
 \E Y_{n,i}^k\leq C_TK^k m^{-ks},\qquad 1\leq i\leq m.
\end{equation}
If $|I_n|<m$, the same estimate follows by increasing $n$: each
piece has at most one step, and $B_n^{-1}=O(2^{-sj_n})$.

Expand $S_n^k$ by the multinomial formula. The pure monomials have
total expectation at most $C_TK^k m^{1-ks}$. In a mixed monomial
$H_n=\prod_iY_{n,i}^{a_i}$, the product runs over the indices with
positive exponents: their sum is $k$, and each is at most $k-1$.
Put $u=k/(k-1)>1$.
Independence, $ua_i\leq k$, and Lyapunov's inequality give
\[
 \E H_n^u=\prod_i\E Y_{n,i}^{ua_i}
 \leq\prod_i(\E Y_{n,i}^k)^{ua_i/k}<C_{m,T,K}.
\]
Thus, for each fixed $m$, every mixed monomial is uniformly
integrable. Also $\sup_n\E X_n^k<\infty$ by
\eqref{tree:eq:internal_diameter_bound}. H\"older's and Markov's
inequalities imply
$\E[H_n\mathbf1_{\{X_n>R\}}]\to0$ uniformly in $n$ as
$R\to\infty$. Since $X_n^k\leq S_n^k$, we obtain
\[
 \lim_{R\to\infty}\limsup_{n\to\infty}
 \E[X_n^k\mathbf1_{\{X_n>R\}}]
 \leq C_TK^k m^{1-ks}.
\]
Only now let $m\to\infty$; $ks>1$ makes the right-hand side vanish.
Finitely many initial $n$ cause no problem, since each has finite
$k$th moment. The convergence in \eqref{tree:eq:moment_convergence}
follows by truncating $x^k$ and then removing the truncation.
Domination proves the uniform-integrability assertion for
smaller costs. Finally, stationarity and
\eqref{tree:eq:internal_norms} give
$\E(X_{I_n}/A_{j_n})^k=1$ when $|I_n|=2^{j_n}$, so
\eqref{tree:eq:moment_convergence} proves the stated normalization
for limits of these internal endpoint costs.
\end{proof}

\subsection{Record scales}

We now choose deterministic scales satisfying the past-ratio bound.
The preceding compactness and moment conclusions will then apply
simultaneously.

\begin{lemma}[Record scales]\label{tree:lem:record_scales}
There are deterministic integers $j_n\to\infty$ and real numbers
$t_n\uparrow p$, $s\leq t_n<p$, such that
\begin{equation}\label{tree:eq:record_ratios}
 \frac{A_{j_n-i}}{A_{j_n}}\leq2^{-t_n i},
 \qquad 0\leq i\leq j_n.
\end{equation}
They may be chosen beyond any prescribed sequence of lower bounds
on $j_n$. Along these scales the moment compactness and uniform
integrability lemmas hold with $K=1$. After extracting so that
every fixed-offset ratio converges, its limit $a_i$ satisfies
\begin{equation}\label{tree:eq:record_exact_power}
 2^{-i}\leq a_i\leq2^{-pi}\quad(i\geq0),\qquad a_0=1.
\end{equation}
\end{lemma}

\begin{proof}
For every $t<p$, Lemma~\ref{tree:lem:moment_growth} implies that
$2^{-tj}A_j$ is unbounded. It therefore has strict record maxima
at arbitrarily large indices. Choose any $t_n\uparrow p$ in
$[s,p)$, and choose $j_n$ to be a strict record for that exponent,
beyond $j_{n-1}$ and any specified lower bound. For $0\leq i\leq
j_n$, the record inequality is
$2^{-t_n(j_n-i)}A_{j_n-i}\leq2^{-t_nj_n}A_{j_n}$, which is
\eqref{tree:eq:record_ratios}. Compactness of each bounded ratio and
a diagonal extraction give \eqref{tree:eq:record_exact_power}.
For its lower bound, iterate \eqref{tree:eq:internal_subadditivity} to get
$A_{j_n}\leq2^iA_{j_n-i}$ before passing to the limit.
\end{proof}

At these records, every weak limit of the normalized unit-interval endpoint
cost has $k$th moment one, and every fixed backward ratio is positive.
The lower bound $2^{-i}$ does not give a uniform positive lower bound
for the reference-unit ratios $2^{ip}a_i$ as $i$ grows. It also gives
no lower bound for the forward ratio $A_{j+1}/A_j$; that is the
normalization issue in the first-failure argument of
Section~\ref{tree:sec:bootstrap}.

We finish by recording the change of units used below. At base index
$j$, one macroscopic unit is $2^j$ walk steps, contour
heights are divided by $2^{j/2}$, and costs by $B_j$. Replacing
the time unit by $v2^j$, $v=2^{-i}$, replaces the height unit by
$v^{1/2}2^{j/2}$. The reference quantum length unit changes by
$v^p$, so costs expressed in reference metric units are divided
by $v^pB_j$. This convention does not assert a scaling law for a
new passage metric. The conditional-law formulation of this
change of units is established in
Proposition~\ref{tree:prop:zoom_transport}.
